\documentclass{scrartcl}

\makeatletter
\DeclareOldFontCommand{\rm}{\normalfont\rmfamily}{\mathrm}
\DeclareOldFontCommand{\sf}{\normalfont\sffamily}{\mathsf}
\DeclareOldFontCommand{\tt}{\normalfont\ttfamily}{\mathtt}
\DeclareOldFontCommand{\bf}{\normalfont\bfseries}{\mathbf}
\DeclareOldFontCommand{\it}{\normalfont\itshape}{\mathit}
\DeclareOldFontCommand{\sl}{\normalfont\slshape}{\@nomath\sl}
\DeclareOldFontCommand{\sc}{\normalfont\scshape}{\@nomath\sc}
\makeatother

\usepackage{xcolor}
\definecolor{blue1}{RGB}{0,102,189}
\definecolor{blue2}{RGB}{98,160,214}
\definecolor{my_orange}{RGB}{243,98,33}

\colorlet{tumblue}  {blue1}
\colorlet{tumorange}{my_orange}

\usepackage{arxiv}
\usepackage{authblk}

\usepackage{url}
\usepackage{graphicx}
\setkeys{Gin}{keepaspectratio} 

\usepackage{tikz}
\usetikzlibrary{mindmap,trees,shadows}
\usetikzlibrary{arrows,intersections}
\usetikzlibrary{angles}
\usetikzlibrary{shapes.misc} \pgfdeclarelayer{background}
\pgfsetlayers{background, main}
\usetikzlibrary{decorations.markings}

\colorlet{incolor}{tumorange}
\colorlet{outcolor}{tumblue}

\usepackage{tikz-dimline}

\tikzset{inNode/.style={draw=incolor, fill=incolor, circle, inner sep=1.5pt}}
\tikzset{outNode/.style={draw=outcolor,  fill=outcolor, circle, inner sep=1pt}}
\tikzset{projNode/.style={draw=outcolor, fill=outcolor, circle, inner sep=1pt}}

\usetikzlibrary{shapes.arrows, fadings, calc, positioning}
\tikzfading[name=fade right,
  left color=transparent!0, right color=transparent!100]
\tikzfading[name=fade left,
  left color=transparent!100, right color=transparent!0]
\tikzfading[name=fade top,
  bottom color=transparent!0, top color=transparent!100]
  
\pgfdeclarelayer{bg}    \pgfdeclarelayer{fg}    \pgfsetlayers{bg,main,fg}  

\definecolor{participantAcolor}{RGB}{243,98,33}
\definecolor{participantBcolor}{RGB}{0,102,189}
\definecolor{participantCcolor}{RGB}{66,70,50}
\definecolor{cplColor}{RGB}{0,0,0}
\definecolor{preciceColor}{RGB}{ 161, 177, 25}

\definecolor{ci_pantone300}{RGB}{ 0, 101, 189}

\definecolor{ci_pantone301}{RGB}{ 0, 82, 147}
\definecolor{ci_pantone540}{RGB}{ 0, 51, 89}

\definecolor{ci_pantone283}{RGB}{ 152, 198, 234}
\definecolor{ci_pantone542}{RGB}{ 100, 160, 200}
\definecolor{ci_ivory}{RGB}{ 218, 215, 203}
\definecolor{ci_orange}{RGB}{ 227, 114, 34}
\definecolor{ci_green}{RGB}{ 162, 173, 0}  
 
% \usepackage{minted}
\usepackage{listings}
\providecommand*{\listingautorefname}{Listing}  \usepackage{amsmath}
\usepackage{amssymb} \usepackage{graphicx}
\usepackage[hidelinks]{hyperref}
\def\chapterautorefname{Chapter}\def\sectionautorefname{Section}\def\subsectionautorefname{Section}\def\subsubsectionautorefname{Section}\def\paragraphautorefname{Paragraph}\def\tableautorefname{Table}\def\equationautorefname{Equation}\pdfstringdefDisableCommands{\def\corref#1{<#1>}}
\pdfsuppresswarningpagegroup=1
\pgfplotsset{compat=1.16}

\usepackage{multirow}
\usepackage{booktabs}
\usepackage[export]{adjustbox}
\usepackage{gensymb}
\usepackage{subcaption}
\usepackage{xspace} \usepackage{algorithm}
\usepackage{algpseudocode}

\usepackage{geometry}
\geometry{a4paper, left=30mm, right=30mm, top=30mm, bottom=30mm}

\newcommand{\incode}[1]{\mintinline{text}{#1}}
\newcommand{\incpp}[1]{\mintinline{c++}{#1}}
\newcommand{\inpython}[1]{\mintinline{python}{#1}}
\newcommand{\infortran}[1]{\mintinline{fortran}{#1}}
\newcommand{\inmatlab}[1]{\mintinline{matlab}{#1}}
\newcommand{\inbash}[1]{\mintinline{bash}{#1}}
\newcommand{\inxml}[1]{\mintinline{xml}{#1}}
\newcommand{\intext}[1]{\mintinline{text}{#1}}

\newcommand\todo[1]{\textcolor{red}{(TODO: #1)}}
\newcommand\comment[1]{\textcolor{gray}{\textit{(#1)}}}
\newcommand\review[1]{\footnote{\textcolor{blue}{\textit{#1}}}}

\usepackage{siunitx}

\usepackage{framed}

\usepackage{enumitem}

\newcommand\github[1]{\href{https://github.com/#1}{\texttt{#1}}}

\setlength{\parindent}{0em}
\setlength{\parskip}{1.2ex}

\hyphenation{preCICE}
\hyphenation{FEniCS}
\hyphenation{OpenFOAM}
\hyphenation{Nutils}

\title{Solving singular IVPs of Lane-Emden type: A numerical comparison of Physics-Informed Neural Network and integration approximation methods}

\author[1,*,$\dagger$]{Carlos Aznarán}
\author[2,*,$\ddagger$]{Matihus Molina}

\affil[1]{Science Department, National University of Engineering}
\affil[2]{Faculty of Mathematical Sciences, National Major University of San Marcos}
\affil[*]{Equal contributors}
\affil[$\dagger$]{Corresponding author. \textit{Email address:} \url{caznaranl@uni.pe} (Carlos Aznarán)}
\affil[$\ddagger$]{Corresponding author. \textit{Email address:} \url{matihus.molina@unmsm.edu.pe} (Matihus Molina)}